% TOP_PROBLEM: {"problemId": "problem.gromov-s-surface-subgroup-question-for-one-ended-hyperbolic-groups", "canonicalTitle": "Gromov's Surface Subgroup Question for One-Ended Hyperbolic Groups", "releaseRank": 323, "primaryDomain": "geometry_topology", "primaryDomainLabel": "Geometry & topology", "displayStatus": "open_with_solved_subcases", "releaseStatus": "open_with_solved_subcases"}
% !TeX program = lualatex
% ATTEMPT_STATUS: unresolved
\documentclass[11pt]{article}
\usepackage[margin=25mm]{geometry}
\usepackage{fontspec}
\setmainfont{Latin Modern Roman}
\usepackage{amsmath,amssymb,mathtools}
\usepackage{unicode-math}
\setmathfont{Latin Modern Math}
\usepackage{xurl,hyperref}
\hypersetup{colorlinks=true,urlcolor=blue}
\setcounter{secnumdepth}{0}
\setlength{\parindent}{0pt}
\setlength{\parskip}{5pt}
\setlength{\emergencystretch}{3em}
\begin{document}
\section*{\texorpdfstring{Gromov's Surface Subgroup Question for One-Ended Hyperbolic Groups}{Gromov's Surface Subgroup Question for One-Ended Hyperbolic Groups}}\label{323}
\subsection{Definitions and mathematical statement}
A finitely generated group is word-hyperbolic if a Cayley graph has uniformly thin geodesic triangles. It is one-ended if complements of sufficiently large finite balls have exactly one unbounded connected component. The surface-subgroup question asks whether every one-ended word-hyperbolic group $G$ contains a subgroup isomorphic to the fundamental group of a closed surface of negative Euler characteristic. One may express the orientable form as
\[
 \exists g\ge2:\quad
 \left\langle a_1,b_1,\ldots,a_g,b_g\ \middle|\
 \prod_{i=1}^g[a_i,b_i]=1\right\rangle\hookrightarrow G.
\]
A closed nonorientable surface of negative Euler characteristic has an orientable double cover, so allowing such surfaces does not change the existence question. A merely immersed surface without injectivity on fundamental groups is insufficient.
\par\smallskip\textit{Formulation and concept references:} \href{https://research-repository.st-andrews.ac.uk/bitstream/handle/10023/26667/Gardam_2022_MRL_SGC_AAM.pdf?sequence=1}{[S1]}.

\subsection{Short English statement}
Does every one-ended hyperbolic group contain the fundamental group of a closed surface of genus at least two?
\subsection{Research attempt: injective surface constructions and the missing global argument}
\textbf{Status and the precise subgroup requirement.} This attempt does not construct a surface subgroup in every one-ended word-hyperbolic group. A surface relation in a presentation, a map of a closed surface into a complex, and an injection of its fundamental group are distinct assertions. The inherited source [S1] discusses the surface subgroup problem. The current source check also found a September 2026 primary manuscript [E1] treating Baumslag doubles along a minimal diskbusting list of words with an explicit total-length bound. Its hypotheses are retained; this is not a theorem about all one-ended hyperbolic groups.

The work below proves finite-index permanence, gives an explicit family of surface doubles, and formulates an injectivity test for further amalgam constructions. These are constructive subcases, with the failure of the general extension stated explicitly.

\textbf{Lemma 323.1 (finite-index permanence).} If $H\leq G$ has finite index, then $G$ contains a closed orientable surface subgroup of genus at least two if and only if $H$ does.

\emph{Proof.} One direction is immediate from inclusion. For the other, suppose $S\leq G$ is isomorphic to $\pi_1(\Sigma_g)$, $g\geq2$. The map of coset sets $S/(S\cap H)\to G/H$ is injective, so $d=[S:S\cap H]$ is finite. The subgroup $S\cap H$ corresponds to a connected $d$-sheeted covering surface of $\Sigma_g$. This cover is closed and orientable, and a lifted cell decomposition shows that Euler characteristic multiplies by $d$. Its genus is therefore
\[
 g'=1+d(g-1)\geq2.
\]
Thus $H$ contains the required subgroup. $\square$

Likewise a closed nonorientable surface group of negative Euler characteristic contains the group of its orientation double cover, which is a closed orientable surface of genus at least two. This explains the equivalence of the two surface conventions in the selected formulation; neither allows a torus or a free group in place of the required subgroup.

\textbf{Why some hypothesis beyond hyperbolicity is necessary.} A nonabelian free group has no closed hyperbolic surface subgroup. Here is the standard covering-space argument with the obstruction made explicit. The free group is the fundamental group of a graph. Every subgroup corresponds to a connected covering graph, whose fundamental group is free: choose a maximal tree and the remaining edges give a free basis. A closed orientable surface group of genus $g\geq2$ is not free. Its closed aspherical surface is a classifying space and has second homology $\mathbb Z$, whereas a graph is a classifying space for a free group and has zero second homology. These homology groups are invariants of the group.

The Cayley graph of a free group of rank at least two is a branching tree. Removing a ball leaves many infinite components, so it is not one-ended. Thus this obstruction is excluded by the actual hypothesis; it does not disprove the selected assertion. It shows why the one-ended condition cannot be silently omitted when reducing to more familiar hyperbolic examples.

\textbf{A complete construction from a boundary word.} Let $g\geq1$ and put
\[
 F=\langle a_1,b_1,\ldots,a_g,b_g\rangle,
 \qquad w=[a_1,b_1]\cdots[a_g,b_g].
\]
This is the fundamental group of a genus-$g$ orientable surface with one boundary component, and $w$ represents the oriented boundary. Take a second copy $F'$ and glue the two boundary circles by an orientation-reversing identification. Van Kampen's theorem gives an amalgam presentation
\[
 F*_{\langle w\rangle=\langle w'\rangle}F',
 \tag{323.1}
\]
with the boundary orientation determining whether the written relation is $w=w'$ or $w=(w')^{-1}$. Reversing the orientation convention on the second surface exchanges the two descriptions.

The resulting space is a closed orientable surface of genus $2g$: the two punctured surfaces have Euler characteristic $1-2g$, gluing along a circle adds no Euler characteristic, and the result has characteristic $2-4g$. Thus (323.1) is itself a closed hyperbolic surface group. The cyclic edge maps are injective because the nonempty reduced word $w$ has infinite order in the free group; its cyclically reduced powers are nonempty. The normal form theorem for free products with amalgamation then also verifies the injections of the two vertex groups.

These groups are one-ended and word-hyperbolic: the closed surface admits a hyperbolic metric, and its universal cover is the hyperbolic plane. A proper cocompact orbit map is a quasi-isometry; the plane is hyperbolic and one-ended. This establishes examples satisfying all the hypotheses, with an explicitly displayed surface subgroup, rather than just examples satisfying the conclusion in an unrelated ambient group.

\textbf{Lemma 323.2 (a checkable amalgam injectivity condition).} Let $G=A*_C B$. Suppose $A_0\leq A$ and $B_0\leq B$ satisfy
\[
 A_0\cap C=C_0=B_0\cap C,
\]
where the intersections use the specified embeddings of the common subgroup $C$ into the two vertex groups. Then the natural homomorphism
\[
 A_0*_{C_0}B_0\longrightarrow A*_C B
\]
is injective.

\emph{Proof.} A nontrivial reduced normal form in the source either lies nontrivially in a vertex group, or is an alternating product of factors from $A_0\setminus C_0$ and $B_0\setminus C_0$, possibly accompanied by a factor in $C_0$. By the intersection assumptions, a factor outside $C_0$ remains outside $C$ in the target. Thus an alternating reduced form remains reduced after applying the map and cannot become the identity by the normal form theorem. A nontrivial element of a vertex group remains nontrivial because the vertex embeddings are injective. $\square$

Consequently, if $A_0$ and $B_0$ are fundamental groups of once-punctured orientable surfaces and $C_0$ identifies their boundary subgroups, the lemma embeds the closed surface obtained by gluing them. A particularly explicit situation is when these punctured-surface free groups sit as free factors of larger free vertex groups and $C$ is exactly the cyclic subgroup generated by the boundary word. Then the required intersections are automatic. To use such an example inside the selected conjecture one must additionally verify that the chosen ambient amalgam is one-ended and hyperbolic; the injectivity lemma itself does not assert those geometric properties.

The intersection hypothesis is not cosmetic. Without it, a syllable outside the source edge subgroup might enter the target edge subgroup, allowing reduction in the target that was impossible in the source. Checking that boundary loops are nontrivial alone would miss this failure.

\textbf{A different sufficient geometric criterion.} Suppose a closed hyperbolic surface admits a local isometry into a locally CAT$(0)$ complex. Passing to universal covers gives a local isometry from the complete hyperbolic plane into a CAT$(0)$ space. The image of any geodesic segment is a local geodesic. In a CAT$(0)$ space every local geodesic is globally geodesic, so its length equals the distance between its endpoints. Thus the lifted map preserves distances and is injective. Equivariance then implies that the original map induces an injection on fundamental groups: an element in the kernel would identify a point with its distinct translate in the universal cover.

This is a useful construction criterion when a compatible local metric is available. An arbitrary finite presentation of a word-hyperbolic group does not provide a CAT$(0)$ presentation complex, and an arbitrary immersed surface need not be a local isometry. Those missing properties cannot be supplied merely by invoking hyperbolicity of the Cayley graph.

\textbf{A surface relation alone gives no injection.} For genus two, send
\[
 a_1\mapsto x,\quad a_2\mapsto y,\quad b_1,b_2\mapsto1
\]
into the free group $\langle x,y\rangle$. The surface relator $[a_1,b_1][a_2,b_2]$ maps to one, so this defines a homomorphism with nonabelian image. It is nevertheless far from injective: both $b_i$ lie in its kernel. More generally, killing all $b_i$ maps a surface group onto a free group. Producing a nontrivial or even nonabelian surface image is therefore not enough.

Conversely, a surjection from a free group onto a surface group does not embed the latter back into the former. A section would imply an embedding, but constructing the section is a new requirement and is impossible in this example by the free-subgroup obstruction above. These two directions are easy to confuse in presentation-based searches.

\textbf{Diagnostics and remaining gap.} The exact finite script verifies the free-word surface relation after the displayed noninjective substitution and checks the covering-genus formula for small degrees. It also reduces powers of the explicit boundary word to confirm the normal-form examples. These calculations test particular words and formulas; they do not decide injectivity for arbitrary surface maps.

The constructions prove the desired subgroup exists for the displayed surface doubles and for amalgams satisfying the stated intersection criterion. The first unproved step toward the general conjecture is obtaining such an injective surface construction, or an equivalent geometric criterion, in an arbitrary one-ended hyperbolic group. The short-word and diskbusting assumptions in the new manuscript [E1] supply extra structure in a specific family and cannot simply be removed. The selected assertion remains unresolved by this attempt.

\subsection{Sources}
\begin{itemize}
\item[S1]Gardam, Kielak, and Logan, The Surface Group Conjectures with Two Generators, Mathematical Research Letters (2022), Section 3.2.
\url{https://research-repository.st-andrews.ac.uk/bitstream/handle/10023/26667/Gardam_2022_MRL_SGC_AAM.pdf?sequence=1}.
\textit{Formulation source.}
\item[E1]Hoang Le Xuan and Nam Hung Tran Nguyen, \emph{Surface subgroups of Baumslag doubles along short words}, arXiv:2609.21191v2, revised 21 September 2026.
\url{https://arxiv.org/abs/2609.21191v2}.
\textit{Primary manuscript consulted 27 September 2026; minimality, diskbusting, and the total-word-length bound are retained in the cited scope.}
\end{itemize}
\end{document}
